% Part: first-order-logic
% Chapter: natural-deduction
% Section: soundness

\documentclass[../../../include/open-logic-section]{subfiles}

\begin{document}

\iftag{FOL}
      {\olfileid[pt]{fol}{ntd}{sou}}
      {\olfileid[pt]{pl}{ntd}{sou}}
      
\olsection{Correção}

\begin{explain}
Um sistema de !!{derivation}, como a dedução natural, é \emph{correto}
se não pode derivar coisas que de fato não seguem.  A correção é
assim uma espécie de propriedade de segurança garantida para sistemas de
!!{derivation}. Dependendo de qual propriedade da Teoria da Prova está em
questão, gostaríamos de saber, por exemplo, que
\begin{enumerate}
\item toda !!{sentence} !!{derivable} é \iftag{FOL}{válida}{uma tautologia};
\item se uma !!{sentence} é !!{derivable} a partir de outras, ela também é uma
  consequência delas;
\item se um conjunto de !!{sentence}s é inconsistente, ele é insatisfatível.
\end{enumerate}
Estas são propriedades importantes de um sistema de !!{derivation}. Se
alguma delas não vale, o sistema de !!{derivation} é deficiente---ele
!!{derive} demais.  Consequentemente, estabelecer a correção de um
sistema de !!{derivation} é da maior importância.
\end{explain}

\begin{thm}[Correção]
\ollabel{thm:soundness}
Se $!A$ é !!{derivable} a partir das hipóteses !!{undischarged}s
$\Gamma$, então $\Gamma \Entails !A$.
\end{thm}

\begin{proof}
Seja $\delta$ uma !!{derivation} de $!A$. Procedemos por
indução no número de inferências em~$\delta$.

Para a base da indução, mostramos a afirmação quando o número de
inferências é~$0$. Neste caso, $\delta$ consiste apenas de uma única
!!{sentence}~$!A$, isto é, uma hipótese. Essa hipótese é
!!{undischarged}, pois hipóteses só podem ser !!{discharged}s por
inferências, e não há inferências.  Assim, qualquer
\iftag{FOL}{!!{structure}~$\Struct{M}$}{!!{valuation}~$\pAssign{v}$}
que satisfaz todas as hipóteses !!{undischarged}s da prova
também satisfaz~$!A$.

Agora o passo indutivo. Suponha que $\delta$ contém~$n$
inferências. A(s) premissa(s) da inferência mais baixa são !!{derive}das
usando sub-!!{derivation}s, cada uma das quais contém menos de~$n$
inferências.  Assumimos a hipótese de indução: as premissas da
inferência mais baixa seguem das hipóteses !!{undischarged}s das
sub-!!{derivation}s que terminam nessas premissas.  Temos que mostrar
que a conclusão~$!A$ segue das hipóteses !!{undischarged}s
de toda a prova.

Distinguimos casos de acordo com o tipo da inferência mais baixa.
Primeiro, consideramos as possíveis inferências com apenas uma premissa.

\begin{enumerate}
\item Suponha que a última inferência é \Intro{\lnot}: a
  !!{derivation} tem a forma
  \begin{prooftree}
    \AxiomC{$\Gamma, \Discharge{!A}{n}$}
    \RightLabel{$\delta_1$}
    \DeduceC{$\lfalse$}
    \DischargeRule{\Intro{\lnot}}{n}
    \UnaryInfC{$\lnot !A$}
  \end{prooftree}
  Pela hipótese de indução, $\lfalse$ segue das hipóteses
  !!{undischarged}s $\Gamma \cup \{!A\}$ de~$\delta_1$. Considere
  \iftag{FOL}{!!a{structure}~$\Struct{M}$}{!!a{valuation}~$\pAssign{v}$}. Precisamos
  mostrar que, se $\iftag{FOL}{\Sat{M}}{\pSat{v}}{\Gamma}$, então
  $\iftag{FOL}{\Sat{M}}{\pSat{v}}{\lnot !A}$. Suponha por redução ao
  absurdo que $\iftag{FOL}{\Sat{M}}{\pSat{v}}{\Gamma}$, mas
  $\iftag{FOL}{\Sat/{M}}{\pSat/{v}}{\lnot !A}$, isto é,
  $\iftag{FOL}{\Sat{M}}{\pSat{v}}{!A}$. Isto significaria que
  $\iftag{FOL}{\Sat{M}}{\pSat{v}}{\Gamma \cup \{!A\}}$. Isto é
  contrário à nossa hipótese de indução. Logo,
  $\iftag{FOL}{\Sat{M}}{\pSat{v}}{\lnot !A}$.

\item A última inferência é \Elim{\land}: há duas variantes: $!A$
  ou $!B$ pode ser inferido da premissa $!A \land !B$. Considere o
  primeiro caso. A !!{derivation}~$\delta$ tem este aspecto:
  \begin{prooftree}
    \AxiomC{$\Gamma$}
    \RightLabel{$\delta_1$}
    \DeduceC{$!A \land !B$}
    \RightLabel{\Elim{\land}}
    \UnaryInfC{$!A$}
  \end{prooftree}
  Pela hipótese de indução, $!A \land !B$ segue das
  hipóteses !!{undischarged}s~$\Gamma$ de~$\delta_1$. Considere
  !!a{structure}~\iftag{FOL}{$\Struct{M}$}{$\pAssign{v}$}. Precisamos
  mostrar que, se $\iftag{FOL}{\Sat{M}}{\pSat{v}}{\Gamma}$, então
  $\iftag{FOL}{\Sat{M}}{\pSat{v}}{!A}$. Suponha
  $\iftag{FOL}{\Sat{M}}{\pSat{v}}{\Gamma}$. Pela nossa hipótese
  de indução ($\Gamma \Entails !A \land !B$), sabemos que
  $\iftag{FOL}{\Sat{M}}{\pSat{v}}{!A \land !B}$. Por definição,
  $\iftag{FOL}{\Sat{M}}{\pSat{v}}{!A \land !B}$ se, e somente se,
  $\iftag{FOL}{\Sat{M}}{\pSat{v}}{!A}$ e
  $\iftag{FOL}{\Sat{M}}{\pSat{v}}{!B}$.  (O caso em que $!B$ é
  inferido de $!A \land !B$ é tratado de forma similar.)
  
\item A última inferência é \Intro{\lor}: há duas variantes: $!A
  \lor !B$ pode ser inferido da premissa~$!A$ ou da
  premissa~$!B$. Considere o primeiro caso. A !!{derivation} tem a forma
  \begin{prooftree}
    \AxiomC{$\Gamma$}
    \RightLabel{$\delta_1$}
    \DeduceC{$!A$}
    \RightLabel{\Intro{\lor}}
    \UnaryInfC{$!A \lor !B$}
  \end{prooftree}
  Pela hipótese de indução, $!A$ segue das hipóteses
  !!{undischarged}s~$\Gamma$ de~$\delta_1$.  Considere
  \iftag{FOL}{!!a{structure}~$\Struct{M}$}{!!a{valuation}~$\pAssign{v}$}. Precisamos
  mostrar que, se $\iftag{FOL}{\Sat{M}}{\pSat{v}}{\Gamma}$, então
  $\iftag{FOL}{\Sat{M}}{\pSat{v}}{!A \lor !B}$. Suponha
  $\iftag{FOL}{\Sat{M}}{\pSat{v}}{\Gamma}$; então
  $\iftag{FOL}{\Sat{M}}{\pSat{v}}{!A}$, pois $\Gamma \Entails !A$ (a
  hipótese de indução). Logo também deve ser o caso que
  $\iftag{FOL}{\Sat{M}}{\pSat{v}}{!A \lor !B}$.  (O caso em que $!A
  \lor !B$ é inferido de~$!B$ é tratado de forma similar.)
  
\item A última inferência é \Intro{\lif}: $!A \lif !B$ é inferido
  de uma subprova com hipótese~$!A$ e conclusão~$!B$, isto é,
  \begin{prooftree}
    \AxiomC{$\Gamma, \Discharge{!A}{n}$}
    \RightLabel{$\delta_1$}
    \DeduceC{$!B$}
    \DischargeRule{\Intro{\lif}}{n}
    \UnaryInfC{$!A \lif !B$}
  \end{prooftree}
  Pela hipótese de indução, $!B$ segue das hipóteses
  !!{undischarged}s de~$\delta_1$, isto é, $\Gamma \cup \{!A\} \Entails
  !B$. Considere
  \iftag{FOL}{!!a{structure}~$\Struct{M}$}{!!a{valuation}~$\pAssign{v}$}. As
  hipóteses !!{undischarged}s de~$\delta$ são apenas $\Gamma$, pois
  $!A$ é descartada na última inferência. Precisamos mostrar que
  $\Gamma \Entails !A \lif !B$.  Por redução ao absurdo, suponha que para alguma
  \iftag{FOL}{!!{structure}~$\Struct{M}$}{!!{valuation}~$\pAssign{v}$},
  $\iftag{FOL}{\Sat{M}}{\pSat{v}}{\Gamma}$ mas
  $\iftag{FOL}{\Sat/{M}}{\pSat/{v}}{!A \lif !B}$. Assim,
  $\iftag{FOL}{\Sat{M}}{\pSat{v}}{!A}$ e
  $\iftag{FOL}{\Sat/{M}}{\pSat/{v}}{!B}$. Mas por hipótese, $!B$ é uma
  consequência de $\Gamma \cup \{!A\}$, isto é,
  $\iftag{FOL}{\Sat{M}}{\pSat{v}}{!B}$, o que é uma contradição. Logo,
  $\Gamma \Entails !A \lif !B$.

\item A última inferência é \FalseInt: aqui, $\delta$ termina em
  \begin{prooftree}
    \AxiomC{$\Gamma$}
    \RightLabel{$\delta_1$}
    \DeduceC{$\lfalse$}
    \RightLabel{\FalseInt}
    \UnaryInfC{$!A$}
  \end{prooftree}
  Pela hipótese de indução, $\Gamma \Entails \lfalse$. Temos que mostrar
  que $\Gamma \Entails !A$. Suponha o contrário; então para
  algum~$\iftag{FOL}{\Struct{M}}{\pAssign{v}}$ temos
    $\iftag{FOL}{\Sat{M}}{\pSat{v}}{\Gamma}$ e
    $\iftag{FOL}{\Sat/{M}}{\pSat/{v}}{!A}$. Mas sempre temos
    $\iftag{FOL}{\Sat/{M}}{\pSat/{v}}{\lfalse}$, de modo que isto
    significaria que $\Gamma \Entails/ \lfalse$, contrariando a hipótese
    de indução.

\item A última inferência é \FalseCl: Exercício.

\iftag{FOL}{
\item A última inferência é \Intro{\lforall}: então $\delta$ tem a forma
  \begin{prooftree}
    \AxiomC{$\Gamma$}
    \RightLabel{$\delta_1$}
    \DeduceC{$!A(a)$}
    \RightLabel{\Intro{\lforall}}
    \UnaryInfC{$\lforall[x][!A(x)]$}
  \end{prooftree}
  A premissa $!A(a)$ é consequência das hipóteses
  !!{undischarged}s $\Gamma$ pela hipótese de indução.  Considere alguma
  estrutura $\Struct{M}$ tal que $\Sat{M}{\Gamma}$.  Precisamos
  mostrar que $\Sat{M}{\lforall[x][!A(x)]}$. Como $\lforall[x][!A(x)]$
  é !!a{sentence}, isto significa que temos que mostrar que para toda
  atribuição de variáveis~$s$, $\Sat{M}{!A(x)}[s]$
  (\olref[syn][ass]{prop:sat-quant}). Como $\Gamma$ consiste inteiramente
  em sentenças, $\Sat{M}{!B}[s]$ para todo $!B \in \Gamma$ por
  \olref[syn][sat]{defn:satisfaction}.  Seja $\Struct{M'}$ como
  $\Struct{M}$, exceto que $\Assign{a}{M'} = s(x)$.  Como $a$ não
  ocorre em~$\Gamma$, $\Sat{M'}{\Gamma}$ por
  \olref[syn][ext]{cor:extensionality-sent}. Como $\Gamma \Entails
  !A(a)$, $\Sat{M'}{!A(a)}$.  Como $!A(a)$ é !!a{sentence},
  $\Sat{M'}{!A(a)}[s]$ por
  \olref[syn][ass]{prop:sentence-sat-true}. $\Sat{M'}{!A(x)}[s]$ se, e somente se,
  $\Sat{M'}{!A(a)}$ por \olref[syn][ext]{prop:ext-formulas} (lembre que
  $!A(a)$ é apenas $\Subst{!A(x)}{a}{x}$). Assim,
  $\Sat{M'}{!A(x)}[s]$. Como $a$ não ocorre em~$!A(x)$, por
  \olref[syn][ext]{prop:extensionality}, $\Sat{M}{!A(x)}[s]$. Mas $s$
  foi uma atribuição de variáveis arbitrária, de modo que
  $\Sat{M}{\lforall[x][!A(x)]}$.
  
\item A última inferência é \Intro{\lexists}: Exercício.

\item A última inferência é \Elim{\forall}: Exercício.
}{}
\end{enumerate}

Agora consideremos as possíveis inferências com várias premissas:
\Elim{\lor}, \Intro{\land}, \iftag{FOL}{\Elim{\lif}, e
  \Elim{\lexists}}{e \Elim{\lif}}.
\begin{enumerate}

\item A última inferência é \Intro{\land}. $!A \land !B$ é inferido
  das premissas $!A$ e $!B$ e $\delta$ tem a forma
  \begin{prooftree}
    \AxiomC{$\Gamma_1$}
    \RightLabel{$\delta_1$}
    \DeduceC{$!A$}
    \AxiomC{$\Gamma_2$}
    \RightLabel{$\delta_2$}
    \DeduceC{$!B$}
    \RightLabel{\Intro{\land}}
    \BinaryInfC{$!A \land !B$}
  \end{prooftree}
  Pela hipótese de indução, $!A$ segue das hipóteses
  !!{undischarged}s~$\Gamma_1$ de~$\delta_1$ e $!B$ segue das
  hipóteses !!{undischarged}s~$\Gamma_2$ de~$\delta_2$. As
  hipóteses !!{undischarged}s de~$\delta$ são $\Gamma_1 \cup
  \Gamma_2$, de modo que temos que mostrar que $\Gamma_1 \cup \Gamma_2 \Entails
  !A \land !B$. Considere
  \iftag{FOL}{!!a{structure}~$\Struct{M}$}{!!a{valuation}~$\pAssign{v}$}
  com $\iftag{FOL}{\Sat{M}}{\pSat{v}}{\Gamma_1 \cup \Gamma_2}$. Como
  $\iftag{FOL}{\Sat{M}}{\pSat{v}}{\Gamma_1}$, deve ser o caso que
  $\iftag{FOL}{\Sat{M}}{\pSat{v}}{!A}$, pois $\Gamma_1 \Entails !A$, e
  como $\iftag{FOL}{\Sat{M}}{\pSat{v}}{\Gamma_2}$,
  $\iftag{FOL}{\Sat{M}}{\pSat{v}}{!B}$, pois $\Gamma_2 \Entails
  !B$. Juntas, $\iftag{FOL}{\Sat{M}}{\pSat{v}}{!A \land !B}$.
  
\item A última inferência é \Elim{\lor}: Exercício.

\item A última inferência é \Elim{\lif}. $!B$ é inferido das
  premissas $!A \lif !B$ e~$!A$. A !!{derivation}~$\delta$ tem este aspecto:
  \begin{prooftree}
    \AxiomC{$\Gamma_1$}
    \RightLabel{$\delta_1$}
    \DeduceC{$!A \lif !B$}
    \AxiomC{$\Gamma_2$}
    \RightLabel{$\delta_2$}
    \DeduceC{$!A$}
    \RightLabel{\Elim{\lif}}
    \BinaryInfC{$!B$}
  \end{prooftree}
  Pela hipótese de indução, $!A \lif !B$ segue das
  hipóteses !!{undischarged}s~$\Gamma_1$ de~$\delta_1$ e $!A$
  segue das hipóteses !!{undischarged}s~$\Gamma_2$
  de~$\delta_2$. Considere
  \iftag{FOL}{!!a{structure}~$\Struct{M}$}{!!a{valuation}~$\pAssign{v}$}. Precisamos
  mostrar que, se $\iftag{FOL}{\Sat{M}}{\pSat{v}}{\Gamma_1 \cup
    \Gamma_2}$, então $\iftag{FOL}{\Sat{M}}{\pSat{v}}{!B}$. Suponha
  $\iftag{FOL}{\Sat{M}}{\pSat{v}}{\Gamma_1 \cup \Gamma_2}$. Como
  $\Gamma_1 \Entails !A \lif !B$, $\iftag{FOL}{\Sat{M}}{\pSat{v}}{!A
    \lif !B}$. Como $\Gamma_2 \Entails !A$, temos
  $\iftag{FOL}{\Sat{M}}{\pSat{v}}{!A}$. Isto significa que
  $\iftag{FOL}{\Sat{M}}{\pSat{v}}{!B}$ (pois se
  $\iftag{FOL}{\Sat/{M}}{\pSat/{v}}{!B}$, como
  $\iftag{FOL}{\Sat{M}}{\pSat{v}}{!A}$, teríamos
  $\iftag{FOL}{\Sat/{M}}{\pSat/{v}}{!A \lif !B}$,
  contradizendo~$\iftag{FOL}{\Sat{M}}{\pSat{v}}{!A \lif !B}$).

\item A última inferência é \Elim{\lnot}: Exercício.
  
\tagitem{FOL}{A última inferência é \Elim{\lexists}: Exercício.}{}
\end{enumerate}
\end{proof}

\tagprob{FOL}
\begin{prob}
Complete a prova de \olref[fol][ntd][sou]{thm:soundness}.
\end{prob}
\tagendprob

\tagprob{notFOL}
\begin{prob}
Complete a prova de \olref[pl][ntd][sou]{thm:soundness}.
\end{prob}
\tagendprob

\begin{cor}
\ollabel{cor:weak-soundness}
Se $\Proves !A$, então $!A$ é \iftag{FOL}{válida}{uma tautologia}.
\end{cor}

\begin{cor}
\ollabel{cor:consistency-soundness}
Se $\Gamma$ é satisfatível, então é consistente.
\end{cor}

\begin{proof}
Provamos a contrapositiva.  Suponha que $\Gamma$ não é consistente.
Então $\Gamma \Proves \lfalse$, isto é, há uma !!{derivation} de
$\lfalse$ a partir de hipóteses !!{undischarged}s em~$\Gamma$. Por
\olref{thm:soundness}, qualquer
\iftag{FOL}{!!{structure}~$\Struct{M}$}{!!{valuation}~$\pAssign{v}$}
que satisfaz $\Gamma$ deve satisfazer $\lfalse$.  Como
$\iftag{FOL}{\Sat/{M}}{\pSat/{v}}{\lfalse}$ para toda
\iftag{FOL}{!!{structure}~$\Struct{M}$}{!!{valuation}~$\pAssign{v}$},
nenhuma \iftag{FOL}{$\Struct{M}$}{$\pAssign{v}$} pode satisfazer $\Gamma$,
isto é, $\Gamma$ não é satisfatível.
\end{proof}

\end{document}
